\documentclass[12pt,a4paper]{article}
\usepackage[utf8]{inputenc}
\usepackage[T2A]{fontenc}
\usepackage[russian]{babel}
\usepackage{amsmath}
\usepackage{amsfonts}
\usepackage{amssymb}
\usepackage{graphicx}
\usepackage{geometry}
\usepackage{indentfirst}
\usepackage{array}
\usepackage{float}
\usepackage{booktabs}
\usepackage{caption}
\usepackage{url}

\geometry{left=2.5cm, right=1.5cm, top=2cm, bottom=2cm}
\graphicspath{{photoes/}}
\newcommand{\COO}{\mathrm{CO_2}}
\title{Лабораторная работа №\,13\\Ответы на контрольные вопросы}
\author{Ларькин Глеб, Б03-301 \\
Хабибуллин Марат, Б03-305 \\
Мурашев Артём, Б03-305 \\
Скворцов Максим, Б03-305}
\date{}
\begin{document}
\maketitle

Вопросы приведены в порядке §~13.6 лабораторного практикума
по физической механике под редакцией Э.~Е.~Сона (МФТИ, 2004).
Для оценок в вопросах 1--3 используем газ при $T\approx300$~К и
$p=1$~атм. При другом давлении характерное время столкновений
меняется примерно обратно пропорционально $p$.

\section{Поступательно-вращательная релаксация}

При нецентральном столкновении сила взаимодействия создаёт момент
относительно центра масс молекулы. Часть поступательной энергии
переходит во вращение. В тепловом газе много доступных вращательных
уровней, поскольку расстояние между соседними заселёнными уровнями
мало по сравнению с $k_BT$. После нескольких столкновений температуры
поступательного и вращательного движения сближаются.

Оценим частоту столкновений по модели твёрдых сфер. Для $\COO$ возьмём
массу молекулы $m\approx7{,}3\cdot10^{-26}$~кг и эффективный
диаметр $d_m\approx4{,}5\cdot10^{-10}$~м:
\[
 n=\frac{p}{k_BT}\approx2{,}4\cdot 10^{25}~\text{м}^{-3},\qquad
 \langle v_{\text{отн}}\rangle=\sqrt{\frac{16k_BT}{\pi m}}
       \approx 540~\text{м/с}.
\]
Частота $\nu\sim n\pi d_m^2\langle v_{\text{отн}}\rangle$
составляет около $8\cdot10^9$~с$^{-1}$, а время между столкновениями
$\tau_c\sim0{,}12$~нс. При ориентировочном числе столкновений
$Z_{RT}\sim2$--$5$ получаем $\tau_{RT}\sim0{,}2$--$0{,}6$~нс.
$Z_{RT}$ зависит от газа и температуры. Такое время существенно
меньше микросекундного времени торможения в опыте, что позволяет
пользоваться общей температурой поступательных и вращательных степеней
свободы.

\section{Вращательно-вращательная релаксация}

При столкновении одинаковых молекул переход одной молекулы вниз по
вращательной лестнице может компенсироваться переходом другой вверх.
Например, обмен состояниями $(J+1,J)\to(J,J+1)$ не требует суммарного
изменения вращательной энергии. При малом энергетическом дефекте
обмен эффективен. Он выравнивает распределение по вращательным уровням,
а общий избыток вращательной энергии уходит в поступательное движение
по механизму из вопроса 1.

Если вероятность такого обмена за столкновение порядка единицы, то
$\tau_{RR}\sim\tau_c\sim10^{-10}$~с. Это оценка масштаба. Для
конкретной пары уровней необходима вероятность перехода с учётом
допустимых изменений углового момента.

\section{Колебательно-колебательная релаксация}

У одинаковых двухатомных молекул обмен первым квантом
$A(1)+A(0)\leftrightarrow A(0)+A(1)$ резонансный. Обмен между
разными парами уровней даёт небольшой дефект из-за ангармоничности,
а между разными газами он зависит от разницы колебательных частот.
Если $|\Delta\omega|t_c\ll1$, где $t_c$~--- длительность одного
столкновения, дефект слабо подавляет вероятность перехода. При $VT$-обмене
в поступательное движение должна перейти энергия целого колебательного
кванта. В теории Ландау--Теллера (§~13.7) возбуждение подавлено при
$\omega t_c\gg1$.

При вероятности обмена $P_{VV}$ среднее время порядка
$\tau_c/P_{VV}$. Например, при $P_{VV}\sim0{,}01$--$0{,}1$ и
$\tau_c\sim0{,}1$~нс получаем $1$--$10$~нс. Числа характеризуют
порядок величины для эффективного обмена. Конкретное значение требует
данных о столкновениях и заселённостях уровней.

\section{Газодинамический лазер на смеси CO$_2$ и N$_2$}

В горячей смеси заселяются колебательные уровни обеих молекул. При быстром
расширении в сопле поступательная температура падает. Уровень $v=1$
азота медленно теряет колебательную энергию и служит её резервуаром.
Его квант ($\sim2331$~см$^{-1}$) близок к энергии асимметричного
колебания $\COO(001)$ ($\sim2349$~см$^{-1}$), поэтому возможна
почти резонансная передача энергии азота углекислому газу.

Нижние лазерные уровни, связанные с симметричной и деформационной модами,
опустошаются быстрее благодаря связи с деформационной модой и её
$VT$-релаксации. Если верхний уровень остаётся более заселённым в расчёте
на одно квантовое состояние, возникает инверсия. В этом случае
излучение на переходах $001\to100$ и $001\to020$
(около 10,6 и 9,6~мкм) усиливается при прохождении через газ.
Различие времён релаксации уровней позволяет получить инверсию после
газодинамического охлаждения.

\section{Измерение по поглощению звука}

Звуковая волна периодически нагревает газ при сжатии и охлаждает при
расширении. Для малых возмущений колебательная температура
$T_v'$ следует уравнению
$\tau_{VT}\,dT_v'/dt+T_v'=T'$. При зависимости $e^{i\omega t}$
получаем $T_v'/T'=1/(1+i\omega\tau_{VT})$. Запаздывание обмена энергией
даёт диссипацию. При малой частоте колебания успевают участвовать в
теплообмене, при большой они практически заморожены. Поглощение за одну
длину волны имеет максимум в области $\omega\tau\sim1$.

В ячейке с газом можно измерить амплитуды на двух расстояниях при
нескольких частотах и получить коэффициент затухания из
$A(x)=A(0)e^{-\alpha x}$. Для определения времени строят $\alpha\lambda$
от частоты и подбирают релаксационную зависимость с учётом вязкого и
теплового фона. При одном процессе грубая оценка
$\tau\sim1/(2\pi f_{\max})$. Например, $6$~мкс соответствует
частоте порядка 25~кГц. Для точного пересчёта акустического времени
в $\tau_{VT}$ учитывают теплоёмкости и связь температур при сжатии.
При фиксированной температуре максимум можно искать и по изменению
давления, поскольку $\tau_{VT}$ примерно обратно пропорционально $p$.

\section{Измерение в ударной трубе}

Ударный фронт быстро повышает поступательную температуру, а колебательная
температура непосредственно за фронтом остаётся близкой к исходной.
Дальше по потоку поступательная энергия переходит в колебания.
По скорости фронта и начальному состоянию из балансов массы, импульса
и энергии определяют параметры сразу за фронтом с учётом замороженных
колебаний.

В зоне релаксации измеряют, например, поглощение ИК-излучения на
подходящей колебательной полосе CO$_2$ или изменение плотности
интерферометром. В системе фронта время вдоль траектории частицы
$t=\int_0^x dx'/u_2(x')$, где $u_2$~--- скорость газа относительно
фронта. При почти постоянном $u_2$ характерная длина зоны
$\ell\sim u_2\tau$. Уравнение энергии и закон релаксации подбирают к
профилю сигнала, учитывая, что температура газа в зоне может меняться.
Так находят характерное время и повторяют измерения для разных
температур за фронтом.

\section{Влияние влажности на релаксацию CO$_2$}

Складываются частоты релаксации при столкновениях с CO$_2$ и с водой.
Если $x$~--- мольная доля воды и $A_i=(p\tau)_i$ для соответствующей
пары столкновений, то при том же полном давлении
\[
 \frac{1}{p\tau_{\text{см}}}=\frac{1-x}{A_{\COO}}+
                  \frac{x}{A_{\mathrm{H_2O}}},\qquad
 \frac{\tau_{\text{см}}}{\tau_{\COO}}
 =\frac{1}{1+x(A_{\COO}/A_{\mathrm{H_2O}}-1)}.
\]
Из рис.~13.1 при комнатной температуре $A_{\COO}\approx6{,}3$~мкс$\cdot$атм,
а в условии $A_{\mathrm{H_2O}}=5\cdot10^{-3}$~мкс$\cdot$атм.
Отношение коэффициентов около 1260. При 100~ppm воды
$\tau_{\text{см}}/\tau_{\COO}\approx0{,}89$, а при 1000~ppm
получается около 0,44. Поэтому даже малая примесь ускоряет релаксацию.

\section{Как меняются теплоёмкость и показатель адиабаты}

Для линейной молекулы поступательное и вращательное движения дают
$3R/2$ и $R$ в $c_v$. При нагреве растёт заселённость возбуждённых
колебательных уровней, и часть подведённого тепла идёт на увеличение
колебательной энергии. Наиболее заметен вклад дважды вырожденной
деформационной моды с наименьшим квантом. Поэтому табл.~13.1 показывает
рост $c_p$. Для идеального газа $c_p-c_v=R$, так что
$\gamma=1+R/c_v$ при этом падает. Предел при замороженных колебаниях
равен $7/5$, а при полном классическом возбуждении четырёх колебательных
степеней свободы $\gamma$ стремится к $15/13$ в модели без диссоциации.

\section{Давление, температура и энтропия по рис.~13.4}

Полная удельная энтальпия $h+V^2/2$ сохраняется в стационарном
адиабатическом течении без внешней работы. В ресивере и после полного
торможения скорость практически нулевая. Поэтому $h_2=h_0$.
Для идеального газа в равновесии $h$ зависит только от $T$, и
$T_2=T_0$. Для реального газа энтальпия может зависеть и от давления,
поэтому равенство температур требует отдельной проверки.

В сопле газ разгоняется и охлаждается приблизительно изэнтропически.
При торможении и последующем выравнивании температур степеней свободы
производится энтропия. При одинаковой конечной и начальной температуре
$S_2-S_0=-R\ln(p_2/p_0)>0$, следовательно, $p_2<p_0$.

Для точки 34 деления при 24\,$^\circ$C берём $T_0=297{,}15$~К,
$T_1\approx281{,}63$~К, $p_0\approx125{,}33$~кПа и $p_1=98{,}66$~кПа.
При $c_p=843$ и $R=188{,}9$~Дж/(кг$\cdot$К) отсчитываем энтропию
от состояния $T_*=273{,}15$~К, $p_*=101{,}325$~кПа:
\[
 S-S_*=c_p\ln(T/T_*)-R\ln(p/p_*).
\]
В ресивере получаем $S_0-S_*\approx71{,}0-40{,}2=30{,}8$~Дж/(кг$\cdot$К),
а на срезе сопла $S_1-S_*\approx25{,}8+5{,}0=30{,}8$~Дж/(кг$\cdot$К).
Их равенство следует из изэнтропической модели разгона.

\section{Расчёт теплоёмкости CO$_2$}

Для гармонического осциллятора средняя тепловая энергия одной моды
на единицу массы равна $R\Theta/(e^{\Theta/T}-1)$, где
$\Theta=hc\tilde\nu/k_B$. Энергия нулевых колебаний постоянна и при
дифференцировании исчезает. Складывая производные по температуре,
получаем
\[
 c_p=\frac72R + R\sum_j g_j
   \frac{(\Theta_j/T)^2 e^{\Theta_j/T}}{(e^{\Theta_j/T}-1)^2}.
\]
Используем волновые числа 1333, 667 и 2349~см$^{-1}$ с кратностями
1, 2 и 1. Они дают характеристические температуры примерно 1918,
960 и 3380~К. Наибольший вклад при 300~К даёт мода 667~см$^{-1}$.
\begin{center}
\begin{tabular}{rcccc}
\toprule
$T$, $^\circ$C & $c_p$ (расчёт) & $c_p$ (табл. 13.1) &
$\gamma$ (расчёт) & $\gamma$ (табл. 13.1) \\
\midrule
50 & 869,5 & 869,0 & 1,2776 & 1,278 \\
75 & 893,7 & 892,7 & 1,2681 & 1,268 \\
\bottomrule
\end{tabular}

\smallskip
$c_p$ в Дж/(кг$\cdot$К).
\end{center}
Расчёт с постоянными частотами согласуется с табличными величинами
с точностью порядка 0,1\,\%.

\section{Вывод формулы (13.23)}

Пусть $x=(T_0-T_1)/T_0$, $c_p=c_{pa}+c_{vib}$ и
$\gamma_V=c_{vib}/c_{pa}$. При мгновенном торможении передача энергии
в колебания за время торможения пренебрежимо мала. Энергия
$V_1^2/2=c_p(T_0-T_1)$ идёт в активные степени свободы:
\[
 T^*=T_1+\frac{V_1^2}{2c_{pa}}=T_0(1+\gamma_V x),\qquad
 p^*=p_1(T^*/T_1)^{c_{pa}/R}.
\]
После торможения колебательная энергия приходит к равновесию при
$p^*$, а температура возвращается к $T_0$ из сохранения полной
энтальпии. Подставляя $p_1=p_0(1-x)^{c_p/R}$, находим
\[
 \Delta S_{\max}=-R\ln\frac{p^*}{p_0}
   =-c_{pa}\ln(1+\gamma_V x)-c_{vib}\ln(1-x).
\]
Линейный член равен $(-c_{pa}\gamma_V+c_{vib})x=0$.
Остаётся квадратичный член
\[
 \Delta S_{\max}\simeq
 \frac{c_{vib}(1+\gamma_V)}{2}
 \left(\frac{T_0-T_1}{T_0}\right)^2,
\]
что и даёт формулу (13.23). Квадратичная зависимость объясняет,
почему при малых скоростях искомый перепад легко сравним с потерями
установки.

\end{document}
